This study investigates whether generating a compressed intermediate representation, later expanded by a deterministic program, yields valid BPMN 2.0 models more reliably and more cheaply than direct XML generation by language models. To that end, the following were built: a domain-specific language with a formal grammar, BPMN-DSL; a transpiler that expands it into XML validated against the official schema; a set of 768 pairs of textual description and DSL; and a seven-billion-parameter model fine-tuned with QLoRA on these pairs, at a cost of US\$~0.47. The evaluation used 46 processes from a public test set, with three generations per item. The transpiler produced valid XML for all 1,021 corpus samples, and the DSL has 81\% fewer tokens than the equivalent XML. With the same large models, asking for the DSL instead of XML lowered validity from 97.8\% to 62.3\% and from 93.5\% to 52.9\%. Most failures were in references to elements already written, which required an identifier declared in advance; with references by label, in an exploratory analysis, validity rose to 87.0\%. The fine-tuned small model also reached 87.0\%, without the grammar in the prompt, with a cost per valid output 5.5 times lower than that of the same method on the large model and 9.4 times lower than that of direct XML generation. By the fidelity measure, which is exploratory because the judge used did not pass the validation criteria, the small model scored about one point lower on a 0 to 5 scale. The conclusion is that the compact notation does not make generation more reliable but makes it cheaper, and that notations meant for autoregressive models should allow references to what has already been written.

\keywords{BPMN; domain-specific language; language models; fine-tuning; token compression.}
